\section{Structural Basis and Polymorphism}

\subsection{Low- and High-Temperature Phases}

BaZn$_2$Si$_2$O$_7$ transforms at roughly 280--300~$^\circ$C. In the low-temperature phase the ZnO$_4$ chains are strongly directional, whereas the high-temperature phase relieves thermal strain by rotating its tetrahedra and adjusting the angles between chains \cite{lin1999phase,thieme2015ba1}. HT\nobreakdash-XRD and dilatometry show that the thermal expansion tensors of the two structures differ by far more than the usual scatter introduced by sintering, which makes the transition itself the primary variable in phase control \cite{kerstan2012thermal,thieme2016very}.

Public structural databases place low-temperature BaZn$_2$Si$_2$O$_7$ in the orthorhombic space group Cmc2$_1$, with Ba$^{2+}$ five-coordinate and both Zn$^{2+}$ and Si$^{4+}$ tetrahedrally coordinated; the ZnO$_4$ and SiO$_4$ tetrahedra share corners to build a three-dimensional network \cite{available2020materials}. These coordination data make a reasonable first-tier screen for structural analogy, but a database entry is not in situ high-temperature structural evidence and is no substitute for the original diffraction refinement. Table~\ref{tab:structure} summarizes the structural information on these phases together with the conclusions that bear on phase control.

\begin{table}[bp]
\centering
\caption{Structure and thermal behavior of BaZn$_2$Si$_2$O$_7$ and its structurally related compounds.}
\label{tab:structure}
\begin{tabular}{P{\dimexpr 0.2102\textwidth-2\tabcolsep\relax}P{\dimexpr 0.2985\textwidth-2\tabcolsep\relax}P{\dimexpr 0.1978\textwidth-2\tabcolsep\relax}P{\dimexpr 0.2935\textwidth-2\tabcolsep\relax}}
\toprule
    \thead{Phase or composition} & \thead{Structure and coordination} & \thead{Characterization} & \thead{Implication for phase control}\\
\midrule
    BaZn$_2$Si$_2$O$_7$ & Orthorhombic; five-coordinate Ba; ZnO$_4$\slash SiO$_4$ tetrahedra & HT\nobreakdash-XRD, dilatometry & Low- to high-temperature transition present \cite{lin1999phase,available2020materials}\\
    \addlinespace[3pt]
    Ba$_{1-x}$Sr$_x$Zn$_2$Si$_2$O$_7$ & Sr on the Ba site; high-temperature structure retained to room temperature & Single-crystal XRD, HT\nobreakdash-XRD & Transition temperature drops sharply for $x\gtrsim0.1$ \cite{thieme2015ba1,thieme2016very}\\
    \addlinespace[3pt]
    BaZn$_{2-x}$Mg$_x$Si$_2$O$_7$ & Mg on the Zn site; structure varies with $x$ & HT\nobreakdash-XRD, dilatometry & Substitution shifts the transition temperature and the expansion anisotropy \cite{kerstan2012thermal}\\
\bottomrule
\end{tabular}
\end{table}

\subsection{Structural Criteria}

The criteria for comparing Ba$_5$Y$_{12}$Zn[O(SiO$_4$)]$_8$ with the reference structure fall into four tiers: whether the silicate groups remain isolated tetrahedra; whether Zn stays four-coordinate and forms traceable corner-sharing chains; whether the Ba\slash Y polyhedra have comparable connectivity dimensionality; and whether the space group and lattice parameters are similar. The first two tiers rest on direct structural evidence for BaZn$_2$Si$_2$O$_7$ \cite{lin1999phase,available2020materials}; the last two require single-crystal or high-resolution powder diffraction data for Ba$_5$Y$_{12}$Zn[O(SiO$_4$)]$_8$, which are not yet available.
